\documentclass[11pt]{article}
\usepackage{amsmath,amsfonts}
\usepackage{graphicx}
%\usepackage{xcolor}
%\definecolor{gray}{rgb}{.75,.75,.75}
%\usepackage[landscape]{geometry}
%\usepackage{hyperref}
%\usepackage[bookmarks=true, pdfborder={0 0 .5}, urlbordercolor=gray]{hyperref}

\renewcommand{\thefootnote}{\fnsymbol{footnote}}


\addtolength{\topmargin}{-.8in}
\addtolength{\textheight}{1.5in}
\addtolength{\oddsidemargin}{-.4in}
\addtolength{\textwidth}{.8in}
\pagestyle{empty}

\newcommand{\ds}{\displaystyle}
\newcommand{\ts}{\textstyle}

\newcommand{\Z}{\mathbb{Z}}
\newcommand{\R}{\mathbb{R}}
\newcommand{\C}{\mathbb{C}}
\newcommand{\EE}{\mathcal{E}}
\newcommand{\tZ}{\tilde{\Z}}

\newcommand{\Sym}{\mathrm{Sym}}

\renewcommand{\circle}[1]{{\bigcirc\hspace{-.75em}#1}}

\begin{document}

\footnote[0]{Notes by Peter Magyar \ {\tt magyar@math.msu.edu}}
\vspace{-1em}

\centerline{\bf Math 133\hspace{1.5em} \hfill{\Large\bf
Method for Convergence Testing}\hfill Stewart \S11.7}
\vspace{1em}

\noindent 
For a series $\sum_{n=1}^\infty a_n=a_1+a_2+a_3+\cdots$, 
determine if it converges toward a limit 
as we add more terms, or diverges (either to $\pm\infty$ or oscillating).
%A key fact is that the first few terms $a_1, a_2,\ldots$
%have no effect on the convergence or divergence, 
%which depends only on the infinite tail of the sum,
%$\sum_{n=N}^\infty a_n$ for any starting point $N$. 

\begin{enumerate}

\item[0.] If $\lim\limits_{n\to\infty} a_n\neq 0$, then the series diverges
by the $n$-th Term Test (Vanishing Test).

\item Try to manipulate the series into a Standard Series:
\begin{itemize}

\item Geometric series: \ $\sum\limits_{n=1}^\infty cr^{n-1} = c+cr+cr^{2} +c r^3+\cdots
=\left\{\!\!\begin{array}{cl}
\frac{c}{1-r} &\text{ for $|r|<1$}\\
\text{diverges} &\text{ for $|r|\geq 1$.}
\end{array}\right.
$
\item Standard $p$-series: \ 
$\sum\limits_{n=1}^\infty \frac{1}{n^p}
= 1 + \frac{1}{2^p} + \frac{1}{3^p}+\cdots
=\left\{\!\!\begin{array}{cl}
\text{converges}  &\text{ for $p>1$}\\
\text{diverges} &\text{ for $p\leq 1$.}
\end{array}\right.
$

\end{itemize}

\item If $a_n$ is a fraction, estimate with a simpler fraction $b_n$, 
often a standard series, by taking only the {\it largest} term from the numerator
and denominator of $a_n$.
\\
Convergence of $\sum a_n$ is usually same as convergence of $\sum b_n$.  
Justify with a Test:
\\[-1.5em]
\begin{itemize}
\item Direct Comparison Test (positive $a_n$)
\begin{itemize}
\item[$\circ$]  Ceiling $0\leq a_n\leq c_n$ where $\sum c_n$ converges \ $\Longrightarrow$ \ $\sum a_n$ also converges.
\item[$\circ$]   Floor $0\leq d_n\leq a_n$ where $\sum d_n$ diverges \ $\Longrightarrow$ \ $\sum a_n$ also diverges.
\end{itemize}
The ceiling $c_n$ or floor $d_n$ will usually be closely related to the estimate $b_n$.
\smallskip
\item Limit Comparison Test: 
Determine $L=\smash{\lim\limits_{n\to\infty}}a_n/\,b_n$. 
\begin{itemize}
\item[$\circ$] $|L|<\infty$ and $\sum b_n$ converges  \ $\Longrightarrow$ \ $\sum a_n$ also converges.

\item[$\circ$]  $|L|>0$ and $\sum b_n$ diverges \ $\Longrightarrow$ \ $\sum a_n$ also diverges. \smallskip

\noindent  [Compare with $(L{-}\epsilon) b_n < a_n< (L{+}\epsilon) b_n$ for $n>N$].\footnote{
Most later tests are proved by reducing to a Direct Comparison, specified in [brackets]}
\end{itemize}
\end{itemize}

\item 
Try the Integral Test if $a_n$ is positive and fairly simple, but not comparable to a standard series: e.g.~$\frac{1}{n\ln(n)}$.
For a positive decreasing function $f(x)$ with $a_n=f(n)$, 
compute  improper integral \smash{$\int_{1}^\infty\!\! f(x)\,dx =\lim\limits_{N\to\infty}\!\!\int_1^N\!\! f(x)\,dx=\lim\limits_{N\to\infty}\!\! F(N){-}F(1)$.}
\begin{itemize}
\item[$\circ$] $\int_1^\infty\!f(x)\,dx$\, converges  \ $\Longrightarrow$ \ $\sum a_n$ also converges \  
[$\sum_{n=1}^\infty a_n\leq a_1+\int_1^\infty\! f(x)\,dx$].
\item[$\circ$] $\int_1^\infty\! f(x)\,dx$\, diverges  \ $\Longrightarrow$ \ $\sum a_n$ also diverges 
[$\sum_{n=1}^\infty a_n\geq \int_1^\infty\! f(x)\,dx$].
\end{itemize}


\item Try the Ratio Test if $a_n$ has a growing number of factors, for example
if it contains $r^n$ or $n!$.
Determine $\smash{\lim\limits_{n\to\infty}} |a_{n+1}/a_n|=L$.
\begin{itemize}
\item[$\circ$] $L<1$  \ $\Longrightarrow$ \ $\sum a_n$ converges \, 
[\,$|a_n|\leq c(L{+}\epsilon)^n$ for $n>N$].
\item[$\circ$] $L>1$  \ $\Longrightarrow$ \  $\sum a_n$ diverges \, 
 [\,$|a_n|\geq c(L{-}\epsilon)^n$ for $n>N$].
\item[$\circ$] $L=1$  \ $\Longrightarrow$ \   no conclusion. [e.g. any standard $p$-series]
\end{itemize}

\item If $\sum a_n$ has positive and negative terms, try:
\begin{itemize}
\item  Absolute Convergence: $\sum |a_n|$ converges \ $\Longrightarrow$ \   $\sum a_n$ also converges. 
\item Alternating Series: $a_n=(-1)^{n-1} b_n$ with $b_n\geq 0$:
\\[.3em]
$\lim\limits_{n\to\infty} b_n=0,\ \text{$b_n$ decreasing} \ \Longrightarrow \  \sum a_n
\text{ converges. }$
\\[.3em]
Error estimate: \ 
$\sum_{n=1}^{2N} a_n\leq \sum_{n=1}^\infty a_n \leq \sum_{n=1}^{2N} a_n + b_{2N+1}$.
\end{itemize}
\end{enumerate}
\end{document}

%%%%%%%%%%%  Picture  %%%%%%%%%%
\\[0em]
\centerline{
%\includegraphics[width=2in]
{1-8-Discont-iv.png}
\qquad \qquad 
%\includegraphics[width=2in]
{1-8-Discont-v.png} 
}
\vspace{0em}

\noindent
%%%%%%%%%%%%%%%%%%%%%%%%%

%%%%%%%%%%%  Table  %%%%%%%%%%
In fact, we get the following derivatives:
$$\begin{array}{|c||c|c|c|c|c|c|}
\hline &&&&&& \\[-.9em] 
f(x) & \sin(x) & \cos(x) & \tan(x)  & \sec(x) & \csc(x) & \cot(x)
\\[.2em] \hline &&&&&& \\[-.9em]
f'(x) &\cos(x) & -\sin(x) & \sec^2(x) & \tan(x)\sec(x) & -\cot(x)\csc(x) & -\csc^2(x)
\\[.1em] \hline
\end{array}$$
\\[1em]
%%%%%%%%%%%%%%%%%%%%%%%%%%



















